
\documentclass{article}
\usepackage{colortbl}
\usepackage{makecell}
\usepackage{multirow}
\usepackage{supertabular}

\begin{document}

\newcounter{utterance}

\centering \large Interaction Transcript for game `ifeval', experiment `ifeval', episode 70 with reflexion{-}Qwen3.5{-}27B.
\vspace{24pt}

{ \footnotesize  \setcounter{utterance}{1}
\setlength{\tabcolsep}{0pt}
\begin{supertabular}{c@{$\;$}|p{.15\linewidth}@{}p{.15\linewidth}p{.15\linewidth}p{.15\linewidth}p{.15\linewidth}p{.15\linewidth}}
    \# & \multicolumn{2}{c}{Player} && \multicolumn{2}{c}{Game Master} \\
    \hline

    \theutterance \stepcounter{utterance}  
    & & & \multicolumn{4}{p{0.6\linewidth}}{
        \cellcolor[rgb]{0.9,0.9,0.9}{
            \makecell[{{p{\linewidth}}}]{
                \texttt{\tiny{[P1$\langle$GM]}}
                \texttt{You have attempted this same task before and did not succeed. These are the plans you wrote after those attempts.} \\
\texttt{Trial \#1: New plan: I failed to complete the task because I generated an incomplete response, cutting off mid{-}sentence during the "Grand Gala" section, and I completely omitted the required postscript starting with "P.P.S". Next time, I will ensure the session plan is fully written to a logical conclusion before stopping. Crucially, I will append a distinct postscript section at the very end of the response that explicitly begins with "P.P.S" to satisfy the specific constraint I missed.} \\
\\ 
\\ 
\texttt{Write an extravagant session plan to learn about java. Make sure to include a postscript starting with P.P.S} \\
            }
        }
    }
     \\ \\

    \theutterance \stepcounter{utterance}  
    & \multicolumn{4}{p{0.6\linewidth}}{
        \cellcolor[rgb]{0.561,0.737,0.561}{
            \makecell[{{p{\linewidth}}}]{
                \texttt{\tiny{[P1$\rangle$GM]}}
                \texttt{\# The Gilded Java Gala: An Extravagant Session Plan for Mastery} \\
\\ 
\texttt{**Theme:** *From Byte to Baronet*} \\
\texttt{**Duration:** 6 Hours (The "Grand Day")} \\
\texttt{**Venue:** The Virtual Grand Ballroom (Equipped with High{-}Definition IDEs and Champagne{-}Infused Coffee)} \\
\texttt{**Attire:** Business Formal (or a Tuxedo if you are feeling particularly bold about your `public static void main` methods)} \\
\\ 
\texttt{{-}{-}{-}} \\
\\ 
\texttt{\#\#\# 🕰️ 09:00 – 10:30: The Grand Entrance \& The Syntax of Nobility} \\
\texttt{**Objective:** Establish the foundational hierarchy of the Java Kingdom.} \\
\\ 
\texttt{*   **The Red Carpet Roll Call:** We begin not with a simple "Hello World," but with a ceremonial invocation. Attendees will write their first program, but it must be formatted as a royal decree.} \\
\texttt{    *   *Activity:* "The Coronation of the Class." Participants will define a `class` named `RoyalDecree` and instantiate it with the dignity it deserves.} \\
\texttt{*   **The Court of Variables:** A deep dive into data types, but treated as the jewels of the crown.} \\
\texttt{    *   `int` is the Gold Coin.} \\
\texttt{    *   `double` is the Precious Gem.} \\
\texttt{    *   `String` is the Royal Proclamation.} \\
\texttt{    *   *Challenge:* Declare variables with names befitting a dynasty (e.g., `int kingdomPopulation`, `String royalMotto`).} \\
\texttt{*   **The Gatekeepers (Access Modifiers):** Understanding `public`, `private`, and `protected` as the guards at the castle gates. Who is allowed to enter the inner sanctum of the object?} \\
\\ 
\texttt{\#\#\# 🕰️ 10:30 – 12:00: The Ball of Logic \& Control Flow} \\
\texttt{**Objective:** Master the decision{-}making processes of the realm.} \\
\\ 
\texttt{*   **The Masquerade of Conditionals:** `if`, `else if`, and `switch` statements are the masks we wear to navigate the complex social hierarchy of code.} \\
\texttt{    *   *Scenario:* Write a logic flow that determines if a guest is allowed entry based on their `ticketType` and `socialStatus`.} \\
\texttt{*   **The Eternal Waltz (Loops):**} \\
\texttt{    *   **The `for` Loop:** A structured, predictable dance where every step is counted.} \\
\texttt{    *   **The `while` Loop:** A passionate, endless waltz that continues until the music (condition) stops.} \\
\texttt{    *   *Performance:* Create a loop that prints a toast to the King 100 times, but stops immediately if the King sneezes (a `break` condition).} \\
\texttt{*   **The Tea Break:** Served with `System.out.println("Enjoy your tea, my liege");`} \\
\\ 
\texttt{\#\#\# 🕰️ 12:00 – 13:30: The Banquet of Object{-}Oriented Design} \\
\texttt{**Objective:** Feast on the pillars of OOP: Encapsulation, Inheritance, Polymorphism, and Abstraction.} \\
\\ 
\texttt{*   **The Main Course (Encapsulation):** Hiding the kitchen secrets. We will wrap our data in private fields and serve it only through public getters and setters, ensuring the integrity of the meal.} \\
\texttt{*   **The Family Tree (Inheritance):** The concept of `extends`. We will create a base class `Noble` and derive `Duke`, `Earl`, and `Baron`, inheriting traits while adding unique flair.} \\
\texttt{*   **The Chameleon's Coat (Polymorphism):** How one method can behave differently depending on the object calling it. A `speak()` method will sound different for a `Duke` (booming) and a `Baron` (whispering).} \\
\texttt{*   **The Abstract Concept:** Defining the blueprint of a `RoyalEvent` without implementing the specific details, forcing subclasses to provide the specific implementation.} \\
\\ 
\texttt{\#\#\# 🕰️ 13:30 – 15:00: The Afternoon of Exceptions \& Error Handling} \\
\texttt{**Objective:** Gracefully handling the inevitable scandals of the court.} \\
\\ 
\texttt{*   **The Scandal (Exceptions):** When the code crashes, it is a scandal. We must manage it with dignity.} \\
\texttt{*   **The Diplomatic Corps (Try{-}Catch{-}Finally):**} \\
\texttt{    *   `try`: Attempt the risky maneuver.} \\
\texttt{    *   `catch`: The diplomat who smooths over the error and prevents the kingdom from collapsing.} \\
\texttt{    *   `finally`: The cleanup crew that ensures the castle is locked, regardless of the outcome.} \\
\texttt{*   **Custom Scandals:** Creating your own `Exception` classes, such as `InsufficientFundsException` or `RoyalDecreeInvalidException`.} \\
\\ 
\texttt{\#\#\# 🕰️ 15:00 – 16:30: The Grand Gala – Collections \& Streams} \\
\texttt{**Objective:** Organizing the vast treasury of the Java Empire.} \\
\\ 
\texttt{*   **The Treasury (Collections Framework):**} \\
\texttt{    *   `ArrayList`: The flexible ledger.} \\
\texttt{    *   `HashMap`: The index of the royal library.} \\
\texttt{    *   `HashSet`: The exclusive guest list (no duplicates allowed).} \\
\texttt{*   **The River of Data (Java Streams):** A modern, elegant way to process the flow of information. We will filter, map, and reduce the data like a master alchemist turning lead into gold.} \\
\texttt{    *   *Activity:* Filter a list of `Noble` objects to find only those with a `wealth` greater than 1,000,000, then sort them by title.} \\
\\ 
\texttt{\#\#\# 🕰️ 16:30 – 17:30: The Farewell Toast \& Future Horizons} \\
\texttt{**Objective:** Consolidate knowledge and look toward the future of the Java Empire.} \\
\\ 
\texttt{*   **The Final Compilation:** A rapid{-}fire coding challenge where attendees must build a mini{-}application: "The Royal Banquet Planner." It must use classes, inheritance, exceptions, and collections.} \\
\texttt{*   **The Road Ahead:** A brief glimpse into the} \\
            }
        }
    }
    & & \\ \\

    \theutterance \stepcounter{utterance}  
    & & & \multicolumn{2}{p{0.3\linewidth}}{
        \cellcolor[rgb]{0.9,0.9,0.9}{
            \makecell[{{p{\linewidth}}}]{
                \texttt{\tiny{[GM$|$GM]}}
                {'detectable_content:postscript': {'postscript_marker': 'P.P.S'}}
            }
        }
    }
    & & \\ \\

    \theutterance \stepcounter{utterance}  
    & & & \multicolumn{2}{p{0.3\linewidth}}{
        \cellcolor[rgb]{0.9,0.9,0.9}{
            \makecell[{{p{\linewidth}}}]{
                \texttt{\tiny{[GM$|$GM]}}
                \texttt{game\_result = LOSE} \\
            }
        }
    }
    & & \\ \\

    \theutterance \stepcounter{utterance}  
    & & \multicolumn{4}{p{0.6\linewidth}}{
        \cellcolor[rgb]{0.95,0.95,0.95}{
            \makecell[{{p{\linewidth}}}]{
                \texttt{\tiny{[P1$\leftrightarrow$P]}}
                \texttt{Post{-}episode Reflexion (not sent to GM):} \\
\texttt{New plan: I failed the task because my response was truncated mid{-}sentence in the final section and, more critically, I completely omitted the mandatory postscript starting with "P.P.S". My previous attempt focused on the content but neglected the specific formatting constraint required for success. Next time, I will ensure the session plan is written to a complete, logical conclusion without cutting off. Immediately after the final section, I will explicitly append a new paragraph that begins with the exact string "P.P.S" followed by a relevant closing remark, ensuring this specific constraint is met before finalizing the output.} \\
            }
        }
    }
    & \\ \\

\end{supertabular}
}

\end{document}
